%% kapa: exam=1 round=22 year=2011 period=1 form=A subject=civil_law subject_ko=민법 seq=6 number=6
%% source: Q-Net 감정평가사 기출문제 (cst00309) artlSeq=5036790 file=extracted/1008/5036790/2011년 제22회 감정평가사 1차시험 기출문제/1교시.hwp
%% conversion: kordoc 4.24.51 -> tools/gen1.py (auto); review=proofread-sonnet; flags=none
\begin{question}{6}
법정추인에 해당하지 않는 경우는?
\begin{choices}
  \item 취소권자가 취소할 수 있는 법률행위로 인하여 발생한 채무의 이행을 상대방에게 청구한 경우
  \item 취소할 수 있는 법률행위로 인하여 발생한 채무를 소멸시키고 그 대신 다른 채무를 성립시키기로 계약한 경우
  \item 취소권자가 상대방으로부터 담보를 제공받은 경우
  \item 취소권자의 상대방이 취소할 수 있는 법률행위로 취득한 권리의 일부를 제3자에게 양도한 경우
  \item 취소권자가 취소할 수 있는 법률행위로 인하여 발생한 채무를 상대방에게 이행한 경우
\end{choices}
\end{question}
\answer{4}
